\begin{examplebox}
	\textbf{\textcolor{CExample}{例 1（基础）}}

	\textbf{\textcolor{CExample}{题目：}}
	三棱锥 $P-ABC$ 中，$PA,PB,PC$ 两两垂直，$PA=3,\ PB=4,\ PC=5$。求该三棱锥外接球的半径。

	\textbf{\textcolor{CExample}{解题步骤：}}
	\begin{description}[style=nextline, leftmargin=2.8em, labelsep=0.5em, itemsep=0.45em]
		\item[\textbf{第一步（识别模型）：}] 由于 $PA,PB,PC$ 两两垂直，符合墙角模型，可补形为长方体。
			\par\textbf{理由：} 墙角模型是补形法的典型应用，三棱锥的三条棱即为长方体的三条棱。
		\item[\textbf{第二步（确定棱长）：}] 设补形所得长方体的三条棱长分别为 $a,b,c$，则 $a=PA=3$，$b=PB=4$，$c=PC=5$。
			\par\textbf{理由：} 墙角模型直接对应长方体共顶点三棱。
		\item[\textbf{第三步（代入公式）：}] 外接球半径 $R = \frac{\sqrt{a^{2}+b^{2}+c^{2}}}{2} = \frac{\sqrt{3^{2}+4^{2}+5^{2}}}{2} = \frac{\sqrt{50}}{2} = \frac{5\sqrt{2}}{2}$。
			\par\textbf{理由：} 补形法得外接球半径公式。
	\end{description}

	\textbf{\textcolor{CExample}{关键结论：}}
	\textbf{墙角模型三棱锥的外接球半径 $R = \frac{\sqrt{a^{2}+b^{2}+c^{2}}}{2}$，其中 $a,b,c$ 为两两垂直的三条棱长。}

	\medskip
	\textbf{\textcolor{CExample}{例 2（稍变形）}}

	\textbf{\textcolor{CExample}{题目：}}
	三棱锥 $A-BCD$ 中，$AB=CD=5$，$AC=BD=\sqrt{34}$，$AD=BC=\sqrt{41}$。求外接球半径。

	\textbf{\textcolor{CExample}{解题步骤：}}
	\begin{description}[style=nextline, leftmargin=2.8em, labelsep=0.5em, itemsep=0.45em]
		\item[\textbf{第一步（判断模型）：}] 三组对棱分别相等，符合对棱相等模型，可补形为长方体。
			\par\textbf{理由：} 对棱相等的三棱锥可补形为长方体，其对棱是长方体相对面的面对角线。
		\item[\textbf{第二步（设未知数）：}] 设长方体共顶点的三条棱长分别为 $x,y,z$，由对应关系得：
			\[
				\begin{cases}
					x^{2}+y^{2} = AB^{2} = 25,\\
					y^{2}+z^{2} = AC^{2} = 34,\\
					z^{2}+x^{2} = AD^{2} = 41.
			\end{cases}
			\]
			\par\textbf{理由：} 长方体相对面的面对角线长度与棱长满足勾股定理。
		\item[\textbf{第三步（解方程组）：}] 三式相加得 $2(x^{2}+y^{2}+z^{2})=25+34+41=100$，故 $x^{2}+y^{2}+z^{2}=50$。分别减各式得 $x^{2}=16,\,y^{2}=9,\,z^{2}=25$，因此 $x=4,\,y=3,\,z=5$。
			\par\textbf{理由：} 代数求解平方和。
		\item[\textbf{第四步（代公式）：}] 补形后的长方体棱长为 $4,3,5$，代入公式 $R = \frac{\sqrt{4^{2}+3^{2}+5^{2}}}{2} = \frac{\sqrt{50}}{2} = \frac{5\sqrt{2}}{2}$。
			\par\textbf{理由：} 直接使用补形法求外接球半径公式。
	\end{description}

	\textbf{\textcolor{CExample}{关键结论：}}
	\textbf{对棱相等的三棱锥，通过解长方体棱长方程组得到 $a,b,c$，再代入公式 $R=\frac{\sqrt{a^{2}+b^{2}+c^{2}}}{2}$ 求解。}

	\medskip
	\textbf{\textcolor{CExample}{例 3（综合应用）}}

	\textbf{\textcolor{CExample}{题目：}}
	四棱锥 $P-ABCD$ 中，$PA\perp$ 底面 $ABCD$，底面 $ABCD$ 是矩形，$PA=4$，$AB=3$，$AD=4$。求四棱锥 $P-ABCD$ 外接球的半径。

	\textbf{\textcolor{CExample}{解题步骤：}}
	\begin{description}[style=nextline, leftmargin=2.8em, labelsep=0.5em, itemsep=0.45em]
		\item[\textbf{第一步（构造长方体）：}] 由 $PA\perp$ 底面 $ABCD$ 且底面为矩形，可将四棱锥补形为长方体：以 $A$ 为原点，$AB,AD,AP$ 为共顶点三条棱。
			\par\textbf{理由：} $PA$ 垂直底面，故可作长方体的高；底面矩形两邻边作为长方体底面两棱。
		\item[\textbf{第二步（确定棱长）：}] 长方体共顶点三棱长分别为 $a=AB=3$，$b=AD=4$，$c=PA=4$。
			\par\textbf{理由：} $AB,AD,AP$ 两两垂直。
		\item[\textbf{第三步（计算半径）：}] 代入公式 $R = \frac{\sqrt{a^{2}+b^{2}+c^{2}}}{2} = \frac{\sqrt{3^{2}+4^{2}+4^{2}}}{2} = \frac{\sqrt{41}}{2}$。
			\par\textbf{理由：} 补形法外接球半径公式。
	\end{description}

	\textbf{\textcolor{CExample}{关键结论：}}
	\textbf{对于有一条侧棱垂直底面且底面为矩形的四棱锥，可补形为长方体，其外接球半径 $R=\frac{\sqrt{a^{2}+b^{2}+c^{2}}}{2}$，$a,b,c$ 分别为矩形相邻边长和侧棱长。}
\end{examplebox}
